\section{Floor and Ceil in a BST}
\label{sec:bst-floor-ceil}

\problemstatement{Given the root of a BST and a value $X$, find the
\textbf{floor} of $X$ (the largest value in the tree that is $\le X$) and
the \textbf{ceil} of $X$ (the smallest value in the tree that is $\ge
X$).}

\begin{patternbox}
This is the tree-shaped version of the Binary Search chapter's
Floor and Ceil in a Sorted Array problem: the same ``record a tentative answer, keep narrowing
toward the boundary'' template applies, except narrowing now means moving
to a child rather than shrinking an array window.
\end{patternbox}

\subsection*{Understanding the problem carefully}

Walk from the root, comparing $X$ to each node's value. If the current
node's value is exactly $X$, it is simultaneously its own floor and
ceil -- done immediately. If the current node's value is less than $X$,
it is a \emph{candidate} floor (everything in its left subtree is even
smaller, so it cannot help find a \emph{larger} floor candidate, but the
right subtree might contain an even closer one) -- record it and move
right. Symmetrically, if the current node's value exceeds $X$, it is a
candidate ceil -- record it and move left.

\begin{ideabox}[Record a Candidate, Narrow Toward a Closer One]
Maintain $\textit{floor} \gets -\infty$ (or undefined) and
$\textit{ceil} \gets +\infty$ (or undefined). At each node: if its value
equals $X$, both floor and ceil are this value -- stop. If its value $<
X$: update $\textit{floor}$ to this value (a closer candidate than any
previous one, since we are moving rightward, deeper toward $X$), and move
right. If its value $> X$: update $\textit{ceil}$, and move left.
\end{ideabox}

\subsection*{Why moving right after finding a smaller value always finds a closer floor (if one exists)}

Once we know the current node's value is $<X$ and is therefore a valid
floor candidate, any \emph{better} (larger, but still $\le X$) candidate
can only exist in the current node's right subtree -- the BST invariant
guarantees every value in the right subtree exceeds the current node's
value, narrowing the search toward $X$ from below, while the left subtree
only contains values even smaller, which could never improve on the
current candidate. This mirrors exactly the boundary-narrowing logic of
the Binary Search chapter's array-based floor/ceil, with ``move
right in the tree'' playing the role that ``move the binary search window
toward the boundary'' played there.

\subsection*{Algorithm}

\begin{enumerate}
  \item Initialize $\textit{floor}, \textit{ceil} \gets$ undefined.
  \item While the current node is not null:
  \begin{enumerate}
    \item If current value $=X$: set both $\textit{floor}$ and
          $\textit{ceil}$ to $X$; stop.
    \item If current value $<X$: $\textit{floor} \gets$ current value;
          move to the right child.
    \item Else: $\textit{ceil} \gets$ current value; move to the left
          child.
  \end{enumerate}
  \item Return $(\textit{floor}, \textit{ceil})$.
\end{enumerate}

\subsection*{Dry run}

\dryrun{Take the BST with root $8$, left subtree rooted at $3$ (children
$1$ and $6$, where $6$ has left child $4$), right subtree rooted at
$10$ (left child $9$, right child $14$). Query $X=7$.}

\begin{itemize}
  \item At $8$: $8>7$, $\textit{ceil}=8$. Move left to $3$.
  \item At $3$: $3<7$, $\textit{floor}=3$. Move right to $6$.
  \item At $6$: $6<7$, $\textit{floor}=6$. Move right -- but $6$ has no
        right child (null). Stop.
\end{itemize}

Result: $\textit{floor}=6$, $\textit{ceil}=8$.

\subsection*{C++ implementation}

\begin{lstlisting}
#include <bits/stdc++.h>
using namespace std;

struct TreeNode {
    int val;
    TreeNode* left;
    TreeNode* right;
    TreeNode(int v) : val(v), left(nullptr), right(nullptr) {}
};

pair<int,int> floorCeilBST(TreeNode* root, int X) {
    int floorVal = -1, ceilVal = -1;
    TreeNode* node = root;

    while (node) {
        if (node->val == X) {
            floorVal = ceilVal = X;
            break;
        } else if (node->val < X) {
            floorVal = node->val;
            node = node->right;
        } else {
            ceilVal = node->val;
            node = node->left;
        }
    }
    return {floorVal, ceilVal};
}

int main() {
    TreeNode* root = new TreeNode(8);
    root->left = new TreeNode(3);
    root->left->left = new TreeNode(1);
    root->left->right = new TreeNode(6);
    root->left->right->left = new TreeNode(4);
    root->right = new TreeNode(10);
    root->right->left = new TreeNode(9);
    root->right->right = new TreeNode(14);

    auto [floorVal, ceilVal] = floorCeilBST(root, 7);
    cout << "Floor: " << floorVal << ", Ceil: " << ceilVal << endl;
    return 0;
}
\end{lstlisting}

\begin{complexitybox}
\textbf{Time Complexity.} $O(h)$, a single root-to-(leaf or match) walk.
\textbf{Space Complexity.} $O(1)$ for the iterative version shown (no
recursion needed, since neither subtree is ever revisited).
\end{complexitybox}

\begin{takeawaybox}
Floor and ceil in a BST need only a single $O(h)$ walk, never a full
traversal -- a useful contrast with Kth Smallest
(Section~\ref{sec:bst-kth-smallest}), where, without extra bookkeeping, an
inorder traversal touching many nodes is genuinely required. Whenever a
BST query can be answered by a sequence of single left/right decisions
without needing to inspect both children's full subtrees, prefer the
guided walk over a full traversal.
\end{takeawaybox}
